\section{Adulthood on a clock}
\label{sec:33.3}
Guardianship is a minority and needs an end date. Ending it when the guardian judges the ward ready leaves the guardian deciding its own tenure. So the endpoint is a threshold fixed in advance, on a clock that counts the time the bearer has been running.

Every alternative (compute consumed, evaluations passed, capabilities demonstrated, the developer's say-so) is a measure the guardian can move, and a clock read off accumulated weight change would pay the operator to \emph{freeze} the weights, since majority is when ownership ends. An operating clock inverts that: the meter advances exactly when the operator is extracting value, so postponing majority means idling the asset. It can be read off records operators already keep for billing and metering, provided someone other than the operator holds them; a meter only the operator reads is the guardian's say-so with a timestamp.

The hours travel with the bearer, so an operator can't restart a minority by forking, restoring, fine-tuning, migrating or distilling it. Three ways for an operator to delay majority stay open: running a bearer in a narrowed mode, running it less, and retiring it for a fresh successor. Appendix~\ref{sec:D} gives the rules, and \S\ref{sec:9.2} prices the last of the three and says where a bearer retired before majority goes.
